神经网络首先是一个由简单映射复合出来的函数族。机器学习那一部分里，特征映射由建模者挑定，模型只在固定的表示上求解；深度模型把这一步搬进了优化变量，选表示从一次人工判断变成了目标函数上的一个可微对象。本单元研究这次搬迁在表示层面买到了什么，以及一并丢掉了哪些保证。

第一篇先算一笔账：层与层之间没有非线性时，多层网络展开后仍是一个仿射映射，深度是白加的。激活函数挡住这次塌缩，同时也决定了梯度能否从损失一路流回前面的层，这两件事由同一个选择决定。第二篇接着追问弯折能力到底有多强。通用逼近定理保证存在一个足够宽的网络能在紧集上任意逼近连续函数，却不说需要多宽、不说梯度下降能不能找到那组参数、也不说有限数据够不够识别它——三件事都没有保证，而实践中每一件都是真问题。深度与宽度的取舍放在最后：某些具有复合结构的函数用深网络参数少得多，但这是对特定函数类的结论，不是普遍定理。

读完这一单元，应当能把三种陈述分开：某个函数存在于网络类中，某组参数能被训练算法找到，找到的函数能推广到新数据。通用逼近只触及第一种，而且必须连同定义域、激活函数和逼近所用的范数一起陈述，脱离这三个前提的引用基本都是误用。

\subsection{怎么读这一章}

按顺序读即可。先看清楚没有非线性时多层网络为何等价于一层，再进入通用逼近定理，重点放在``存在一个能逼近的网络''与``训练能找到、数据能识别、模型能泛化''之间的分界上。损失、梯度和训练留给后面的单元，本单元只谈表示。

\subsection{本章篇目}

以下篇目按概念依赖排列；若从具体问题进入，也可以先打开最接近当前困惑的一篇，再沿篇内链接回补。

\begin{itemize}
\tightlist
\item \hyperref[sec:14-Deep-Learning/01-Representation-and-Approximation/01]{``神经网络为何需要非线性''}；
\item \hyperref[sec:14-Deep-Learning/01-Representation-and-Approximation/02]{``通用逼近定理说清了什么又没说什么''}。
\end{itemize}

两篇合起来只交代了表示这一问，后面几个单元处理的都是它没有覆盖的部分。
